\section{An exact algebraic solution at the first later truncation}
\label{sec:exact-M2}

The sharp-Euler incoming report explicitly proposed exact elimination
and boundary comparison for $K_2$ as a first target
\cite[Section 8]{SharpEulerIncoming}. Both can be completed with small
polynomials. This gives an exact answer, including uniqueness, rather
than an enclosure of a numerically located stationary point.

\begin{theorem}[The unique maximum of $K_2$]\label{thm:M2}
Let $y_2$ be the unique root in $(0,1)$ of
\begin{equation}\label{eq:M2-R}
 R(y)=y^7+3y^6-2y^5-10y^4+17y^3+59y^2-4.
\end{equation}
Define
\begin{align}
 D(y)&=y^6+2y^5-2y^4-6y^3+17y^2+36y+4,
 \label{eq:M2-D}\\
 r_2&=12/D(y_2),\qquad M_2=K_2(r_2,y_2).
 \label{eq:M2-r}
\end{align}
Then $(r_2,y_2)$ is the unique global maximizer of $K_2$ on
$[0,1]^2$. The number $M_2$ has degree seven over $\mathbb Q$ and
satisfies
\begin{equation}\label{eq:M2-minpoly}
\begin{split}
 0={}&32M_2^7-416M_2^6+2538M_2^5-4523M_2^4\\
    &-32216M_2^3-66352M_2^2+99072M_2+20736.
\end{split}
\end{equation}
In particular $M_2=1.11893868596105\ldots$.
\end{theorem}
\begin{proof}
The kernel simplifies to
\begin{equation}\label{eq:K2-rational}
 K_2(r,y)=\frac{r(1+y)(3-r-ry^2-r^2y^2)}{(1+r)(1+ry^2)}.
\end{equation}
Its partial derivatives, after removing positive denominators and
the negative factors $-(1+y)$ and $-r$, have numerators
\begin{align*}
 P(r,y)={}&r^4y^4+2r^3y^4+2r^3y^2+r^2y^4+8r^2y^2\\
          &+r^2+2ry^2+2r-3,\\
 Q(r,y)={}&r^3y^4+r^2y^4+2r^2y^2+6ry^2+8ry+r-3.
\end{align*}
The exact polynomial identities used in the elimination are
\begin{equation}\label{eq:M2-resultant}
 \operatorname{Res}_r(P,Q)=-192y^{10}(y+1)R(y),
\end{equation}
and the penultimate subresultant
\begin{equation}\label{eq:M2-subresultant}
 16y^{10}\bigl(rD(y)-12\bigr).
\end{equation}
These are finite algebraic identities; the accompanying verifier
reconstructs them from the derivative polynomials.

There is exactly one zero of $R$ in $(0,1)$. Indeed $R(0)=-4$,
$R(1)=64$, and
\[
 R'(y)=y\bigl(7y^5+18y^4-10y^3-40y^2+51y+118\bigr)>68y.
\]
For each such $y$, $P(r,y)$ is strictly increasing in $r\ge0$,
and $P(0,y)=-3$, $P(1,y)=4y^4+12y^2>0$. Thus at most one
interior stationary point can occur.

We next prove that a global maximum is interior. At $r=0$ the
kernel is zero, and at $y=0$ it is $1-W_2(r)\le1$. At $y=1$,
\[
 K_2(r,1)=\frac{2r(1-r)(3+r)}{(1+r)^2}<1,
\]
because the denominator minus the numerator equals
$(1-2r)^2+r^2+2r^3>0$. At $r=1$ and $0<y<1$,
$\partial_pK_2<0$ by $P(1,y)>0$, so moving into the square
increases the value. Finally
\[
 K_2(3/4,1/4)=\frac{8325}{7504}>1.
\]
Compactness therefore supplies an interior global maximum.
It must be the unique stationary point identified above.
The subresultant gives $r_2=12/D(y_2)$; its denominator is positive,
since $D(y)\ge y^6+2y^5+9y^2+36y+4$ on $(0,1)$.

Substituting $r=12/D(y)$ in \eqref{eq:K2-rational}, clearing its
positive denominator, and eliminating $y$ against $R(y)$ gives
\eqref{eq:M2-minpoly}. To establish the degree assertion, reduce its
primitive polynomial modulo $5$ and make it monic. The exact certificate
checks
\[
 \gcd(f,X^5-X)=1,\qquad X^{5^7}\equiv X\pmod f.
\]
Since $7$ is prime, these are the finite-field irreducibility criterion
in degree seven. Its leading coefficient does not vanish modulo $5$,
so Gauss's lemma proves irreducibility over $\mathbb Q$.
Rational bisection of $R$ followed by interval evaluation of
\eqref{eq:M2-r} certifies the displayed decimal prefix.
\end{proof}

The geometric significance of this computation is limited but precise.
It solves a pointwise maximization problem for the remainder kernel.
It does not identify the order-averaged $C_2$: the two gamma shape
parameters cannot be assumed to concentrate independently at
$(r_2,y_2)$.

